\chapter{أفراد فريدون}\label{ch:g9:unique-individuals}

سبعة مليارات إنسان، ولا اثنان متشابهان --- لا في عائلة،
ولا في ساحة مدرسة، ولا في التاريخ كله، والتوأمان
المتطابقان مستثنيان نصف استثناء. وقد جمعت ثلاثة فصول
الآن آلة تلك الحقيقة كلها. وهذا الفصل الجامع القصير يجري
الحجة مرة واحدة من طرفها إلى طرفها، ويودع نتائجها: في
تنوع \omterm{def:g5:biodiversity:species}{النوع}، وفي كيف نقرأ العناوين ``\omterm{def:g9:heredity-and-traits:heredity}{الوراثية}''، وفي ما
تعنيه الفرادة وما لا تعنيه.

\section{مبرهنة الفرادة}

\begin{proposition}[لماذا لا يتطابق شخصان]\label{prop:g9:unique-individuals:theorem}
اجمع الآليات الثلاث:
\begin{enumerate}
  \item \emph{خلط التنصيف}: \omterm{def:g8:sexual-reproduction:gametes}{أمشاج} كل والد تحمل شريكًا
        واحدًا من كل زوج، موزعًا مستقلًا --- $2^{23}$
        نصف مجموعة لكل والد
        (\cref{rem:g9:chromosomes:whyhalve})، بل أكثر من
        ذلك بكثير في الحقيقة، إذ \emph{يتبادل الشريكان
        قطعًا} أيضًا في أثناء الانقسام المنصّف، فتقطع
        المجموعة دون مستوى \omterm{def:g9:chromosomes:chromosomes}{الصبغيات} الكاملة؛
  \item \emph{يانصيب \omterm{def:g5:flower-to-fruit:steps}{الإخصاب}}: فأي \omterm{def:g8:reproductive-systems:male}{نطفة} من الملايين
        (\cref{ex:g8:reproductive-systems:sperm}) تلقى
        \omterm{def:g4:animal-reproduction:sexual}{بويضة} أي شهر أمر يضرب الفضاءين أحدهما في الآخر؛
  \item \emph{رشح الطفرات}: فكل جيل يضيف بضع تهجئات
        جديدة من عنده
        (\cref{prop:g9:genes-and-dna:explains}).
\end{enumerate}
وفضاء التركيبات يقزّم عدد البشر الذين عاشوا قط: فكل
\omterm{def:g4:animal-reproduction:sexual}{إخصاب} يوزع مجموعة لم توزع من قبل، ولن توزع من بعد.
\end{proposition}
\admitted

\begin{example}[التوأمان، وحدود الاستثناء]\label{ex:g9:unique-individuals:twins}
التوأمان المتطابقان يتقاسمان توزيعًا واحدًا
(\cref{exo:g8:from-egg-to-birth:10}) --- بل هما نفسهما
يفترقان: فالحياتان المنفصلتان تكتبان \omterm{def:g9:heredity-and-traits:traits}{صفات} مكتسبة مختلفة
(\cref{prop:g9:heredity-and-traits:sorting})، وتواريخ
\omterm{def:g8:nervous-communication:synapse}{مشابك} مختلفة
(\cref{ex:g8:nervous-communication:juggler})، بل بضع
طفرات متأخرة مختلفة في \omterm{def:g5:first-look-at-cells:cell}{خلايا} كل جسم. فالشق الوحيد في
الفرادة يلتئم من نفسه: بعد التوزيع، يفرّق العيش.
\end{example}

\begin{remark}[فريد، وشبيه في أكثره]\label{rem:g9:unique-individuals:alike}
أمسك بالحقيقتين معًا: نصا \omterm{def:g9:genes-and-dna:dna}{DNA} عند أي إنسانين يتفقان في
أكثر من $99$ في المئة --- وهو البناء المشترك \omterm{def:g5:biodiversity:species}{لنوع}
\cref{prop:g5:first-look-at-cells:principle} الواحد ---
والاختلافات، مبثوثة في ثلاثة مليارات حرف، تكفي مع ذلك
لفرادة صارمة. ``مختلفون كلهم'' و``أقرباء كلهم'' حقيقة
واحدة بتكبيرين
(\cref{met:g9:genes-and-dna:levels})؛ وكل استعمال
للحقيقة الأولى ضد الثانية كان بيولوجيا رديئة قبل أن
يكون أخلاقًا رديئة.
\end{remark}

\section{الفرادة في العمل}

\begin{example}[التعرف بواسطة DNA]\label{ex:g9:unique-individuals:forensics}
الفرادة قابلة للفحص، والمحاكم تفحصها: فقراءة ما يكفي من
مواضع الشخص المتغيرة، من أثر \omterm{def:g5:first-look-at-cells:cell}{خلايا}، تطابق فردًا واحدًا
(والتوأمان المتطابقان جانبًا) من البشرية كلها. والمقارنة
نفسها، إن أجريت بتكبير العائلة، تحسم النسب --- إذ يجب أن
يوزع كل موضع متغير عند الولد من مجموعتي والديه. فمنطق
الألبوم (\cref{ch:g9:heredity-and-traits}) قد صار أداة.
\end{example}

\begin{example}[نقل الدم وزرع الأعضاء]\label{ex:g9:unique-individuals:medicine}
الطب يفاوض الفرادة كل يوم. فزمر \omterm{prop:g4:journey-of-food:absorb}{الدم}
(\cref{ex:g9:genes-and-dna:bloodgroups}) هي الحالة
السهلة --- أربع فئات عريضة، تطابق من السجل. أما زرع
الأعضاء فيلقى الحالة الصعبة: إذ يقرأ الجهاز المناعي
(\cref{ch:g9:immune-defenses} يشرح آلته) علامات ذات أدق
حبيبة، فريدة إلى حد يجعل المتبرعين يبحث عنهم في العائلة
أولًا --- فكلما قربت \omterm{prop:g6:classification-kinship:kinship}{القرابة} قرب التوزيع.
\end{example}

\begin{proposition}[التنوع احتياطي النوع، مودعًا]\label{prop:g9:unique-individuals:reserve}
وعد \Cref{rem:g8:sexual-reproduction:reserve} بغاية
الخلط؛ والآليات تصرفها الآن. \omterm{def:g5:biodiversity:species}{فالنوع} المؤلف من أفراد
فريدين يحمل، مبثوثة بينهم، \omterm{def:g9:genes-and-dna:gene}{أليلات} لكل مناسبة ---
المقاومة التي يحملها بعضهم لهذا المرض، والتحمل الذي
يحمله آخرون لذلك الإجهاد (بستان
\cref{ex:g8:sexual-reproduction:bets}، على مقياس
\omterm{def:g5:biodiversity:species}{النوع}). وحين تنقلب الظروف، تقع كوابح
\cref{prop:g8:reproduction-and-environment:limits} على
غير سواء عبر ذلك التنوع --- وما يفعله ذلك التفاوت،
جيلًا بعد جيل، هو بعينه موضوع الفصلين التاليين: تغيّر
\omterm{def:g5:biodiversity:species}{الأنواع}.
\end{proposition}
\admitted

\begin{method}[تدقيق ادعاء ``وراثي'']\label{met:g9:unique-individuals:claims}
في أي عنوان عن المورّثات والناس:
\begin{enumerate}
  \item فحص المستوى
        (\cref{met:g9:genes-and-dna:levels}): أهو ادعاء
        عن مورّثة، أم عن مجموعة، أم عن نمط عائلي؟
  \item وفحص النسيج
        (\cref{prop:g9:heredity-and-traits:sorting}):
        أيحترم نسيج المدى الموروث المعيش فيه، أم يتظاهر
        بأن مورّثة واحدة تعني مصيرًا واحدًا؟
  \item وفحص التنوع: أيعامل انتشار \omterm{def:g8:reproduction-and-environment:population}{الجمهرة} بوصفه
        الاحتياطي الذي هو، أم بوصفه انحرافات عن مجموعة
        ``عادية'' لا وجود لها؟
  \item وفحص التوأمين: أيصمد أمام التوأمين المتطابقين
        اللذين يتقاسمان التوزيع ويختلفان مع ذلك؟
\end{enumerate}
\end{method}

\section{تمارين}

\begin{exercise}[$\star$]\label{exo:g9:unique-individuals:1}
سمّ آليات مبرهنة الفرادة الثلاث.
\end{exercise}

\begin{exercise}[$\star$]\label{exo:g9:unique-individuals:2}
ماذا يضيف تبادل القطع في أثناء التنصيف إلى الخلط؟
\end{exercise}

\begin{exercise}[$\star$]\label{exo:g9:unique-individuals:3}
اذكر شطري \cref{rem:g9:unique-individuals:alike} ---
النسبة المئوية والفرادة --- في جملة واحدة.
\end{exercise}

\begin{exercise}[$\star$]\label{exo:g9:unique-individuals:4}
كيف ينتهي التوأمان المتطابقان إلى أن يميز أحدهما من الآخر؟
\end{exercise}

\begin{exercise}[$\star$]\label{exo:g9:unique-individuals:5}
عن أي سؤالين تجيب مقارنة \omterm{def:g9:genes-and-dna:dna}{DNA} أمام المحكمة؟
\end{exercise}

\begin{exercise}[$\star$]\label{exo:g9:unique-individuals:6}
لماذا يبحث عن متبرعي الزرع في العائلة أولًا؟
\end{exercise}

\begin{exercise}[$\star\star$]\label{exo:g9:unique-individuals:7}
اضرب الحجة: اجمع فضاءي نصفي مجموعتي الوالدين ويانصيب
\omterm{def:g5:flower-to-fruit:steps}{الإخصاب} في نتيجة ``لم توزع من قبل''.
\end{exercise}

\begin{exercise}[$\star\star$]\label{exo:g9:unique-individuals:8}
اصرف \cref{prop:g9:unique-individuals:reserve} على كرم
\cref{exo:g8:sexual-reproduction:10}: ما الذي افتقرت
إليه المستنسخات وتملكه \omterm{def:g8:reproduction-and-environment:population}{جمهرة} نامية من \omterm{def:g2:from-seed-to-plant:seed}{البذور}؟
\end{exercise}

\begin{exercise}[$\star\star$]\label{exo:g9:unique-individuals:9}
طبّق \cref{met:g9:unique-individuals:claims} على:
``اكتشاف مورّثة للنجاح المدرسي''.
\end{exercise}

\begin{exercise}[$\star\star$]\label{exo:g9:unique-individuals:10}
لماذا كانت عبارة ``المجموعة البشرية العادية'' بلا مرجع؟
أجب بالاحتياطي والمبرهنة.
\end{exercise}

\begin{exercise}[$\star\star$]\label{exo:g9:unique-individuals:11}
يبرئ فحص نسب أبًا شكّ فيه جار
\cref{exo:g9:heredity-and-traits:11} ظلمًا. فسّر ما الذي
يفحصه الفحص، موضعًا موضعًا.
\end{exercise}

\begin{exercise}[$\star\star\star$]\label{exo:g9:unique-individuals:12}
``مختلفون كلهم، أقرباء كلهم، وبالآلية نفسها.'' اكتب
الفقرة الجامعة: الخلط، \omterm{prop:g9:genes-and-dna:explains}{والطفرة}، والشفرة المشتركة، كل
واحد في موضعه.
\end{exercise}

\section{مسألة: حلقة القضية القديمة}

\begin{problem}[{مسألة نهاية الأسبوع --- شعرة واحدة،
وثلاثة أسئلة، وثلاثون سنة}]\label{pb:g9:unique-individuals:1}
حلقة دراسية (متخيَّلة) عن قضية قديمة: سرقة لم تحل منذ
ثلاثين سنة، وشعرة واحدة محفوظة \omterm{def:g5:first-look-at-cells:cell}{بخلايا} جذرها، وطرائق
مقارنة حديثة. وتسير الحلقة بالطلبة عبر ما يستطيع التعرف
بواسطة \omterm{def:g9:genes-and-dna:dna}{DNA} أن يقوله وما لا يستطيع.

\medskip\noindent\textbf{الجزء الأول --- ما تحمله
الشعرة.}
\begin{enumerate}
  \item لماذا وجب حفظ \emph{\omterm{def:g1:plants-around-us:parts}{جذر}} الشعرة لقراءة كاملة
        (\cref{def:g6:cell-unit-of-life:parts} يعرف أين
        يسكن \omterm{def:g9:genes-and-dna:dna}{DNA})؟
  \item يقرأ المختبر طقمًا من المواضع الشديدة التغير.
        ولماذا المتغيرة --- ماذا كانت المواضع المشتركة
        بنسبة 99 في المئة ستثبت؟
  \item البصمة تطابق مشتبهًا به في كل موضع. اذكر ما
        يؤكده التطابق، واستثناءه المنهجي الوحيد
        (\cref{ex:g9:unique-individuals:twins}).
  \item يلاحظ الدفاع أن للمشتبه به توأمًا متطابقًا في
        الخارج. فماذا على المحكمة أن تفعل الآن، ولماذا؟
\end{enumerate}

\medskip\noindent\textbf{الجزء الثاني --- أسئلة
العائلة.}
\begin{enumerate}[resume]
  \item بصمة سابقة جزئية كانت قد أشارت إلى ``قريب من
        عائلة ك''. فسّر كيف يقرأ التطابق \emph{الجزئي}
        \omterm{prop:g6:classification-kinship:kinship}{قرابة} (توزيع
        \cref{prop:g9:unique-individuals:theorem}، مجرى
        إلى الوراء).
  \item يسأل محلّف هل تكشف البصمة \omterm{def:g9:heredity-and-traits:traits}{صفات} المشتبه به ---
        المظهر، والمزاج، و``مورّثات الإجرام''. طبّق
        \cref{met:g9:unique-individuals:claims} على
        السؤال، فحصًا فحصًا.
  \item ويسأل محلّف آخر هل يمكن أن يكون التطابق
        ``واحدًا في هذه المدينة'' لا واحدًا في البشرية.
        ما الذي يحدد عدد المواضع المقروءة؟
  \item ويحذر قائد الحلقة: ``البيولوجيا تقول \omterm{def:g5:first-look-at-cells:cell}{خلايا} من
        هي؛ وعلى القضية أن تقول كيف وصلت إلى هناك.''
        فرّق بين الادعاءين.
\end{enumerate}

\medskip\noindent\textbf{الجزء الثالث --- ختام
الحلقة.}
\begin{enumerate}[resume]
  \item لخّص الأداة في جملة واحدة: أي آليتي فصلين
        (\cref{ch:g9:chromosomes}،
        \cref{ch:g9:genes-and-dna}) تجعلان البصمة فريدة
        والعائلة مقروءة.
  \item ولماذا تعمل الأداة نفسها على شعرة عمرها ثلاثون
        سنة؟ (أي خاصية في \omterm{def:g9:genes-and-dna:dna}{DNA} --- بمعزل عن النسخ في
        \cref{prop:g9:genes-and-dna:explains} --- يعوَّل
        عليها: الثبات.)
  \item وتنتهي الحلقة على الأخلاق: اذكر استعمالين
        للفرادة يبينهما هذا الفصل في خدمة الناس، وسوء
        الاستعمال التاريخي الذي يحذر منه
        \cref{rem:g9:unique-individuals:alike}.
  \item اختم بالمبرهنة معادة الصياغة لقاعة المحكمة: ما
        الذي لم يحدث مرتين قط، ولماذا.
\end{enumerate}
\end{problem}
